\subsection{The Original Term Domain and the Remaining Inverse-Basic Progress Step}
\label{sec:p12-equational-basic-completeness}

The earlier definitions matter to the scope of source Theorem~4.13.
They permit open terms and open target substitutions. Initial image
normalization can be justified without introducing new constants;
the separate obstacle is progress from R-equivalent open normal
endpoints. The results below isolate that obstacle and give an actual
inverse-basic trace that preserves normal witness images while
increasing the source's first complexity coordinate.
\cite[\S2.1, pp.~451--452; \S4.3, pp.~495--498]{BaaderSnyder2001}

\begin{hbdefinition}{def:p12-eb-source-scope}{The source's term, substitution and convergence domains}
In Baader--Snyder \S2.1, \(\mathcal F\) is finite or countably
infinite, \(\mathcal V\) is countably infinite, and terms belong to
\(T(\mathcal F,\mathcal V)\). Nullary symbols belong to
\(\mathcal F\); no infinite reserve of nullary symbols is required.
Substitutions map variables to those terms and have finite support.
Their composition acts left to right. Definition~3.1 permits arbitrary
such substitutions as equational unifiers, not only ground ones.
Definition~4.3 calls a substitution R-reduced when all its nonidentity
images are R-normal. A B-state solution R-unifies P but literally
unifies C,S. These are different equality requirements.
\cite[pp.~451--452,469--470,489,495]{BaaderSnyder2001}

The rewriting terminology explicitly imported by \S4.1 is defined
on all finite terms over function symbols, constants and variables.
A rewrite step includes every contextual substitution instance of
an R-rule; it has no distinct-endpoint proviso. A derivation may
therefore contain a step whose endpoints coincide. A reduction
ordering is well founded, monotone and fully substitution-invariant.
Ordinary convergence means termination and confluence on the entire
term domain; ground confluence restricts the latter condition to
variable-free terms of that vocabulary. The two definitions are
not identified. Baader--Snyder's \S4.3/Theorem4.13 expressly says
ground convergent and does not separately declare closure under
free-constant extensions.
\cite[\S2, pp.~541--543; Definitions4.3--4.5, pp.~546--547;
Definitions5.5--5.6, pp.~560--561]{DershowitzPlaisted2001Rewriting}
\cite[pp.~488,496--497]{BaaderSnyder2001}
\end{hbdefinition}

\begin{hbdefinition}{def:p12-eb-ground-map}{One-ground-term evaluation is a proof map, not an infinite-support source substitution}
When \(q\in T(\mathcal F,\varnothing)\), let \(H_q\) evaluate
every variable at q and preserve every function symbol. The map acts
on individual finite terms. It need not be a finite-support
substitution on the infinite alphabet. On each finite set of terms
it agrees with the finite substitution assigning their finitely
many occurring variables to q. On any one rewrite step, a finite
restriction to the chosen rule variables realizes its grounded
instance. This distinction also applies when an infinite sequence
uses infinitely many variable names overall.
\end{hbdefinition}

\begin{hblemma}{thm:p12-eb-ground-termination}{proved=annex:p12.eb.ground-termination}{Ground termination supplies open termination when one ground term exists}
Under the source's unfiltered rewrite relation, if the original
signature has at least one ground term and R terminates on all
its ground terms, then R terminates on all its open terms.
This requires neither distinct ground codes for parameters nor
new constants. It implies no open confluence. If equal-endpoint
rewrite instances were discarded, or if no ground term existed,
the stated implication would need additional hypotheses.
\end{hblemma}

\begin{hblemma}{thm:p12-eb-initial-normalization}{proved=annex:p12.eb.initial-normalization}{Initial open reduced witnesses exist at the actual ordering scope}
Suppose \(l\succ r\) for every R-rule under a reduction ordering
on the source term domain. Then R terminates on open terms;
each left side is nonvariable and
\(\operatorname{Var}(r)\subseteq\operatorname{Var}(l)\).
For every finite-support R-unifier \(\theta\) of an initial
\(P;\varnothing;\varnothing\), normalize its finitely many
nonidentity images by any finite R-reductions. The resulting
finite-support \(\theta_0\) is R-reduced, still R-unifies P,
and is pointwise R-equal to \(\theta\). Thus \(\theta_0\) and
\(\theta\) cover each other on the initial problem variables
with the identity factor. No idempotence of the target substitution,
grounding, fresh-constant reserve or uniqueness of image normal
forms is needed. This is normalization of an initial witness while
S is empty, not alteration of an active literal binding.
\end{hblemma}

\begin{hbtheorem}{thm:p12-eb-open-history-lift}{proved=annex:p12.eb.open-history-lift}{Open normal images, actual helper factors and protected basic history}
Consider the six primitive B transformations, with the source solved
state and actual U helper. Keep a finite-support witness \(\rho\)
R-reduced, R-solving P and literally solving C,S. Literal constraint
steps preserve that witness, including its globally retained initial
scope. If a term of \(P\rho\) has a chosen innermost forward R-redex,
its address is nonvariable in P. A fresh rule match gives normal
proper-subterm images, and the actual matching-first U block succeeds
with computed \(\eta\) and explicit factor
\(\eta\rho'=\rho'\), where \(\rho'\) preserves all old values.
The new solved substitution is \(\sigma_S\eta\), and
\(\sigma_S\eta\rho'=\rho'\) globally.

P's eligible function occurrences originate in the initial skeleton
or a fresh LP right-side template, never in Basic Elimination.
Each protected substitution image has a normal value under the
maintained witness; hence a selected forward redex cannot be inside
it. These assertions hold for open witnesses and need no common-
normal-form or ground-confluence assumption. They do not assert
that a forward redex exists whenever two open endpoints are R-equal.
Nor does forward normality rule out inverse expansions inside an
image. The latter are excluded by the basic addresses, not by a
false bidirectional-normality claim.
\end{hbtheorem}

\begin{hblemma}{thm:p12-eb-inverse-expansion}{proved=annex:p12.eb.inverse-expansion}{A normal eligible inverse match has a compatible normal-witness extension}
Under the same ordering and primitive-state standing, let u be a
nonvariable basic skeleton occurrence with \(u\rho\) R-normal.
Suppose a fresh reversed R-rule \(r\approx l\), from forward
\(l\to r\), has an actual match \(r\gamma=u\rho\).
Its matched variable images are normal; any variables in l absent
from r can be assigned fresh normal parameter variables. The
matching-first helper then supplies a finite raw inverse-basic
block preserving all old witness values and producing the actual
replacement \(l\gamma\). All new function occurrences come from
the eligible fresh l-template; protected substituted images stay
normal. This replacement is strictly greater than \(u\rho\)
under \(\succ\), so it is not a descent step for the source M.
The lemma proves the chosen local inverse operation, not existence
of a suitable eligible inverse match in every unsolved state.
\end{hblemma}

\begin{hbtheorem}{thm:p12-eb-progress-countercase}{proved=annex:p12.eb.progress-countercase}{A source-scope progress failure with a successful raw inverse-basic branch}
Over \(\mathcal F=\{a,f,g\}\), where a is nullary and f,g
unary, take
\[
 R=\{f(x)\to g(x),\quad f(g(x))\to g(x),\quad g(a)\to a\}.
\]
R respects a ground-total reduction ordering, is left-linear and regular, and
terminates on all open terms. It is ground confluent in this
original signature: every closed term reduces to a. Nevertheless
the open terms \(g(g(y))\) and \(g(y)\) are distinct R-normal
terms and are R-equal. Identity is an R-reduced solution of
\(P=\{g(g(y))=^{?}g(y)\}\), but neither endpoint has a
forward redex. Any R-reduced substitution R-equal to identity
on \(\{y\}\) also maps y to y, so choosing a different normal
image representative does not repair the missing forward step.

In the raw inverse-identity B reading, choose fresh
\(g(v)\approx f(v)\) at the left root. Its matching-first
constraint gives \(v\approx g(y)\), with a normal image.
The retained skeleton becomes \(f(v)=^{?}g(y)\).
Choose the forward rule \(f(g(w))\to g(w)\), fresh w.
First matching Decomposition gives \(g(w)=^{?}g(y)\), then
literal solving gives \(w\approx y\). Constrain on the remaining
\(g(w)=^{?}g(y)\) deletes a literal identity. The terminal
S is \(\{v\approx g(y),w\approx y\}\), so its answer on y
is identity. Normal witness values and basic addresses are preserved
throughout. The first inverse step strictly increases the source
multiset M: \(f(g(y))\succ g(g(y))\). Thus the source's forward-
descent cases do not by themselves exhaust this actual term scope,
even though the raw inverse-basic branch succeeds.
\end{hbtheorem}

\begin{proofboundary}{def:p12-eb-progress-boundary}{What remains after initial normalization and the local history proof}
Initial finite-image normalization is supplied above under the
source's actual reduction-ordering hypothesis. It is not the
unresolved dependency. At the literal fixed-signature ground-
convergence scope, what remains is a global progress argument for
R-equal but different open normal endpoints: for example, a finite
equational-conversion-to-basic construction with a justified measure
and occurrence transport for inverse expansions. The source's
grounding step does not provide that argument without establishing
the requisite convergence in the grounding vocabulary.

The countercase refutes the exhaustion/descent mechanism at that
scope, not raw inverse-basic completeness itself. Its successful
inverse trace prevents transferring a directed failure to the
larger family. Strengthening confluence to the whole open domain
would make a different theorem; a ground-target-only result would
also be a different statement. Neither is silently substituted for
the original source claim. Existing stronger-domain completeness,
terminal normalization and arbitrary-theory hereditary refinements
retain their own proof boundaries.
\cite[Theorem4.13 proof, pp.~497--498]{BaaderSnyder2001}
\end{proofboundary}
